\documentclass[12pt,a4paper]{article}

\usepackage[utf8]{inputenc}
\usepackage[T2A]{fontenc}
\usepackage[russian]{babel}
\usepackage[margin=2.5cm]{geometry}
\usepackage{parskip}
\usepackage{longtable}
\usepackage{array}

\title{Translation \\}
\author{}
\date{}

\begin{document}
\maketitle

\section*{Original}

Some previous literature has applied PH to image segmentation, but
the PH calculation has typically been applied to the input image and
used as a way to generate features which can then be used by another
algorithm. Applications have included tumour segmentation [17], cell
segmentation [18] and cardiac segmentation from computed tomography
(CT) imaging [19].

The important distinction between these methods and our approach is
that we apply PH not to the input image being segmented, but rather
to the candidate segmentation provided by the network. To the best
of our knowledge, the first work to take this approach was our
preliminary work in [4], although this idea has subsequently been
developed for the specific case of one-dimensional topological
features (i.e., connected components) in [20]. Here, we extend our
preliminary work [4] by introducing an explicit topological loss
function which can be used to introduce prior knowledge of any
topological feature(s). We also include more extensive experiments
on two different medical imaging modalities (including one
three-dimensional modality) as well as the MNIST dataset.

By applying PH to the candidate segmentations of a neural network,
the topological information found by the PH calculation can be used
to provide a training signal to the network, allowing us to compare
the topological features present in a proposed segmentation with
those specified to exist by some prior knowledge. Importantly, this
can be done even if those topological features are not easily
extracted from pixel intensities in the original image. The
mathematical details that describe how PH quantifies the presence of
topological features in an image, or a candidate segmentation, are
introduced in the next section.

Persistent homology is a method for calculating the robustness of topological features of a dataset at different scales. It
is part of an emerging field known as topological data analysis, 
in which ideas from topology are used to extract information from noisy and high-dimensional datasets. PH has
found applications in neuroscience [21], studying phase
transitions [22], analysis of tree-structured data [23], and in
measuring image artefacts [24]. Below we give an overview
of PH as it applies to our method, but for more thorough
reviews we direct the reader to [11], [25], [26].
PH is most often applied to data forming a high-dimensional point cloud, and the topology of simplicial complexes
generated from that point cloud is the object of study. In our
case though, the data derives from 2D images or 3D volumes, and so a cubical complex is a more natural representation.
 A cubical complex is a set of points, unit line
segments, unit squares, cubes, hypercubes, etc. We define
an elementary interval as a closed subset of the real line 

\section*{Перевод}

В ряде предыдущих работ персистентные гомологии (ПГ) уже
применялись к сегментации изображений однако расчёт ПГ как
правило выполнялся для входного изображения и использовался как
способ генерации признаков которые затем применялись каким-либо
другим алгоритмом. Среди областей применения сегментация
опухолей сегментация клетоки сегментация сердца по
данным компьютерной томографии (КТ).

Важное отличие этих методов от нашего подхода состоит в том, что мы
применяем ПГ не к входному изображению подлежащему сегментации а
к кандидатной сегментации полученной от сети. Насколько нам
известно, первой работой использующей такой подход, была наша
предварительная работа хотя впоследствии эта идея была развита
для частного случая одномерных топологических признаков (т.е.
связных компонент). В данной работе мы расширяем нашу
предварительную работу вводя явную топологическую функцию
потерь которая может использоваться для внедрения априорных знаний
о любых топологических признаках. Мы также включаем более обширные
эксперименты на двух разных модальностях медицинской визуализации
(включая одну трёхмерную модальность) а также на наборе данных
MNIST.

Применяя ПГ к кандидатным сегментациям, полученным нейронной сетью
мы можем использовать топологическую информацию найденную в
результате расчёта ПГ в качестве обучающего сигнала длч сети что
позволяет сравнивать топологические признаки присутствующие в
предложенной сегментации с признаками заданными априорным знанием.
Важно, что это можно сделать, даже если эти топологические признаки
нелегко извлечь из значений интенсивности пикселей в исходном
изображении. Математические детали описывающие то как ПГ
количественно оценивает наличие топологических признаков на
изображении или в кандидатной сегментации представлены в следующем
разделе.

Персистентная гомология - это метод вычисления устойчивости
топологических особенностей набора данных в различных масштабах.
Он является частью развивающейся области, известной как топологический анализ данных, в которой идеи из топологии 
используются для извлечения информации из зашумленных и многомерных наборов данных. Персистентная гомология нашла применение в нейробиологии, 
при изучении фазовых переходов,  анализе древовидных структур данных и при измерении артефактов изображения. 
Ниже мы приводим обзор персистентной гомологии в применении к нашему методу, но для более подробного ознакомления рекомендуем обратиться к источникам.

Персистентная гомология чаще всего применяется к данным, образующим многомерное облако точек, и объектом изучения является топология симплициальных
комплексов, построенных на основе этого облака точек. Однако в нашем случае данные получены из 2D изображений или 3D объектов, поэтому более естественным 
представлением является кубический комплекс. Кубический комплекс - это набор точек, единичных отрезков, единичных квадратов, кубов, гиперкубов и так далее.
Элементарным интервалом мы называем замкнутое подмножество вещественной прямой.

\newpage
\section*{Glossary}

\begin{longtable}{p{3.2cm}p{5.3cm}p{6.8cm}}
\textbf{Term} & \textbf{English explanation} & \textbf{Russian explanation} \\
\hline
\endhead
persistent homology (PH) &
A method from topological data analysis that tracks the birth and
death of topological features (components, loops, cavities) across
a range of scales/thresholds. &
Персистентные гомологии - метод топологического анализа данных,
отслеживающий <<рождение>> и <<смерть>> топологических признаков
(компонент, петель, полостей) на разных масштабах/порогах. \\

topology &
The branch of mathematics studying properties of shapes that are
preserved under continuous deformation, such as connectivity. &
Топология - раздел математики, изучающий свойства форм,
сохраняющиеся при непрерывной деформации, например связность. \\

candidate segmentation &
The segmentation output currently produced by the network, prior to
being evaluated or corrected. &
Кандидатная сегментация - текущий результат сегментации,
выданный сетью, до его оценки или коррекции. \\

topological feature &
A property such as a connected component, a loop/hole, or a cavity
that characterises the shape of an object. &
Топологический признак - свойство (например, связная компонента,
петля/отверстие или полость), характеризующее форму объекта. \\

connected component &
One of the separate, internally connected pieces that make up a
segmented object. &
Связная компонента - одна из отдельных, внутренне связанных
частей, из которых состоит сегментируемый объект. \\

one-dimensional topological feature &
Here: a topological feature of dimension zero in the Betti-number
sense used casually to mean connected components (as opposed to
loops or cavities). &
Здесь: топологический признак, соответствующий связным компонентам
(в отличие от петель или полостей). \\

topological loss function &
A differentiable loss term that penalises deviation of a
segmentation's topology from prior knowledge. &
Топологическая функция потерь - дифференцируемая составляющая
функции потерь, штрафующая отклонение топологии сегментации от
априорного знания. \\

prior knowledge &
Information about the object (here, its expected topology) known in
advance, independent of the specific image. &
Априорное знание - информация об объекте (здесь - его ожидаемая
топология), известная заранее, независимо от конкретного
изображения. \\

training signal &
Information derived from a loss function that is used to update a
network's weights via gradient descent. &
Обучающий сигнал - информация, получаемая из функции потерь и
используемая для обновления весов сети методом градиентного спуска. \\

pixel intensity &
The numerical value representing brightness/colour of a pixel in an
image. &
Интенсивность пикселя - числовое значение, отражающее яркость
или цвет пикселя изображения. \\
\end{longtable}

\section*{Words}

\begin{longtable}{p{2.6cm}p{5.4cm}p{7.3cm}}
\textbf{Word} & \textbf{English meaning} & \textbf{Russian meaning} \\
\hline
\endhead
calculation &
The process of computing PH for a given input. &
Расчёт, вычисление (здесь - вычисление персистентных гомологий). \\

feature &
A distinguishing characteristic extracted from data (here,
topological). &
Признак - отличительная характеристика, извлекаемая из данных
(здесь - топологическая). \\

distinction &
A difference that sets one thing apart from another. &
Отличие, различие - то, что отделяет одно от другого. \\

preliminary &
Coming before the main or final work; introductory. &
Предварительный - предшествующий основной или итоговой работе. \\

subsequently &
Afterwards, at a later point. &
Впоследствии, позднее. \\

extensive &
Covering a large area or range; thorough. &
Обширный, широкий - охватывающий большой объём или диапазон. \\

modality &
A particular type or method of medical imaging (e.g., MRI, CT,
ultrasound). &
Модальность - конкретный вид или метод медицинской визуализации
(например, МРТ, КТ, УЗИ). \\

quantify &
To express or measure something as a numerical value. &
Количественно оценивать, измерять числом. \\
\end{longtable}

\section*{Phrases and Collocations}

\begin{longtable}{p{4.3cm}p{5cm}p{5.9cm}}
\textbf{Phrase / Collocation} & \textbf{English explanation} & \textbf{Russian explanation} \\
\hline
\endhead
to apply PH to X &
To use the persistent-homology calculation on object X (an image or
a segmentation). &
Применять ПГ к X - использовать расчёт персистентных гомологий
для объекта X (изображения или сегментации). \\

to the best of our knowledge &
A hedging phrase used to state that, as far as the authors are
aware, a claim is true (common in academic writing). &
Насколько нам известно - оговорка, используемая для утверждения,
верного в пределах осведомлённости авторов (типично для научного
стиля). \\

to extend preliminary work &
To build upon and develop further an earlier, less complete study. &
Расширять предварительную работу - развивать и дополнять более
раннее, менее полное исследование. \\

to introduce a loss function &
To define and add a new loss function to a training procedure. &
Вводить функцию потерь - определять и добавлять новую функцию
потерь в процедуру обучения. \\

to provide a training signal &
To supply information (via gradients) that guides the updating of a
network's weights. &
Обеспечивать обучающий сигнал - предоставлять информацию (через
градиенты), направляющую обновление весов сети. \\

to compare X with Y &
To examine two things (here, predicted vs. prior-specified
topological features) in order to note similarities/differences. &
Сравнивать X с Y - сопоставлять два объекта (здесь - предсказанные
и заданные априори топологические признаки) для выявления сходств
или различий. \\

to extract features from pixel intensities &
To derive characteristics of an image directly from its raw pixel
values. &
Извлекать признаки из значений интенсивности пикселей - получать
характеристики изображения непосредственно из его исходных
значений. \\

the mathematical details are introduced in the next section &
A standard academic transition phrase pointing the reader to the
following section for formal/technical content. &
Математические детали представлены в следующем разделе -
стандартная переходная фраза научного текста, отсылающая читателя к
формальному/техническому содержанию далее. \\
\end{longtable}

\end{document}
